\chapter*{Ringraziamenti}
% Spazio sotto il titolo: prima erano due righe vuote (~\newline), che davano avvisi
% "Underfull"; questo valore riproduce la stessa posizione della prima riga.
\vspace*{\dimexpr 3\baselineskip - 3.2pt\relax}
% Pagine senza numero ne' testata, come frontespizio e dedica: la numerazione romana
% parte dagli abstract (main.tex azzera il contatore subito prima). Con lo stile
% "fancy" del template, un chapter* stamperebbe inoltre la testata "Capitolo 0".
% \chapter* impone da solo lo stile "plain" alla propria prima pagina: va annullato.
\pagestyle{empty}
\thispagestyle{empty}

\setlength{\parskip}{1em}

% Testo da guideline/dediche.docx, riportato parola per parola.

Ci sono traguardi che, per quanto personali, non si raggiungono mai da soli. Per questo desidero dedicare alcune righe a tutte le persone che, in modi diversi, hanno camminato al mio fianco in questi anni, condividendo con me fatiche, soddisfazioni e momenti indimenticabili.

Desidero innanzitutto rivolgere il mio più sincero ringraziamento alla Prof.ssa Paola Viterbori, relatrice di questa tesi e già relatrice del mio elaborato di laurea triennale. Grazie per aver continuato ad accompagnarmi anche in questo percorso, guidandomi con competenza, disponibilità e grande attenzione in ogni sua fase. La Sua presenza costante, anche nei momenti più intensi del lavoro, la cura con cui ha seguito ogni dettaglio, la gentilezza e la generosità nel condividere la Sua esperienza hanno rappresentato per me un prezioso punto di riferimento. Aver avuto l'opportunità di essere seguita da Lei sia durante il percorso triennale sia in quello magistrale è stato un privilegio e una preziosa occasione di crescita personale e professionale, che porterò sempre con me con profonda stima e sincera ammirazione.

Un ringraziamento va anche alla Prof.ssa Maria Carmen Usai, correlatrice di questa tesi, per la disponibilità al confronto e l'attenzione dedicata a questo lavoro.

Desidero ringraziare di cuore Carlotta Rivella, che mi ha accompagnata durante questo progetto di ricerca. Grazie per la tua costante disponibilità, per il tempo che mi hai dedicato e per i preziosi suggerimenti che hanno accompagnato le diverse fasi del lavoro. Il confronto con te, i tuoi consigli e il tuo supporto sono stati per me importanti e mi hanno aiutata ad affrontare la ricerca con maggiore serenità e consapevolezza.

Un sincero ringraziamento va agli asili nido che hanno aderito al progetto di ricerca e a tutte le educatrici che ci hanno accolte con disponibilità, aprendoci le porte delle loro sezioni e permettendoci di entrare nelle loro routine quotidiane, spesso intense e frenetiche. La fiducia e la collaborazione che ci avete dimostrato hanno reso possibile questo lavoro. Un pensiero va anche alle bambine e ai bambini che hanno partecipato alla ricerca: con la loro spontaneità, curiosità e naturalezza hanno dato ancora più significato a questa esperienza, ricordandomi il valore della ricerca educativa.

A Roberta Tripodi, Giorgia Spitaleri e Ludovica Schiano, compagne di questo progetto, grazie per l'impegno, la dedizione e la collaborazione che hanno caratterizzato ogni fase del nostro lavoro. Condividere questa esperienza con voi ha significato poter contare su un confronto costante e su un supporto reciproco, affrontando insieme tanto le fasi più semplici quanto quelle più impegnative.

A Roberta Tripodi, Giorgia Spitaleri e Costanza Nuti, grazie per aver reso questi due anni di magistrale molto più di due anni di università. Grazie per i caffè prima delle lezioni, per le pause tra una lezione e l'altra, per le giornate trascorse insieme, per le risate, per le ansie condivise prima degli esami e per la capacità di rendere più leggeri anche i momenti più impegnativi. In questi anni siete diventate molto più che compagne di corso: siete diventate amiche. Grazie per il sostegno, dentro e fuori dall'università, e per aver condiviso con me non solo questo percorso accademico, ma anche tanti momenti di vita. Mi auguro che ciò che abbiamo costruito in questi anni possa continuare ben oltre l'università.

Ai miei genitori, grazie. Siete stati e continuate a essere il porto sicuro a cui so di poter sempre tornare, le persone che più di chiunque altro hanno gioito per ogni mio traguardo e che non hanno mai smesso di essermi accanto.

Grazie per tutti i sacrifici che avete affrontato lungo il mio percorso di studi, senza mai farmeli pesare e permettendomi di inseguire i miei obiettivi con serenità. Mi avete sempre lasciata libera di scegliere la mia strada, incoraggiandomi in ogni decisione e facendomi sentire che, qualunque scelta avessi fatto, avrei sempre potuto contare su di voi.

Siete stati presenti nei giorni più belli, ma ancora di più in quelli più difficili, quelli in cui riesco a mostrare la parte più fragile e, ammettiamolo, anche quella più antipatica di me. Eppure non mi avete mai fatto mancare la vostra vicinanza.

Papà, grazie per esserci sempre stato, con una presenza che è andata ben oltre il sostegno emotivo. Grazie per tutto il tempo che hai trascorso con me a Pietra Ligure, rendendo più semplice la mia vita da fuori sede con i tuoi passaggi, la spesa e tutti quei gesti concreti che, uno dopo l'altro, mi hanno permesso di dedicarmi con maggiore serenità allo studio. La tua disponibilità e la naturalezza con cui ti sei sempre preso cura di me hanno reso questo percorso molto più sereno.

Mamma, grazie per aver ascoltato infinite volte le ripetizioni degli esami che più temevo, anche quando probabilmente avresti preferito parlare di qualsiasi altra cosa. Grazie per aver ascoltato argomenti che ormai, dopo tutti questi anni, potresti quasi sostenere al posto mio e, soprattutto, per aver saputo ascoltare me: le mie preoccupazioni, le mie insicurezze e tutte quelle volte in cui avevo semplicemente bisogno che qualcuno mi ricordasse che ce l'avrei fatta.

Ad Anna Garbellini e Nicola Lorenzini, grazie per essere, da sempre, i miei ``genitori putativi''. Gli impegni, gli studi e la quotidianità non ci permettono sempre di viverci quanto vorremmo e, a volte, passa più tempo di quanto desidereremmo tra un incontro e l'altro. Questo, però, non ha mai cambiato la certezza di sapervi presenti nella mia vita. Grazie per l'affetto e il sostegno che non mi avete mai fatto mancare e per essere, da sempre, due persone su cui so di poter contare.

Un grazie va a tutta la mia famiglia, che ha sempre fatto il tifo per me e ha condiviso ogni piccolo e grande traguardo.

Nonno, grazie per esserci sempre stato con il tuo affetto e il tuo sostegno. Ogni volta che mi sentivi dire: ``Nonni, vado a studiare'', la tua risposta era quasi sempre la stessa: ``Ancora?''. E quando ti spiegavo che avevo un esame, concludevi sorridendo: ``Gli esami non finiscono mai''. Dietro a quelle parole ho sempre percepito la tua curiosità nel seguire il mio percorso universitario e il tuo orgoglio nel vedermi crescere e raggiungere, passo dopo passo, i miei obiettivi. Sono piccoli momenti che porterò sempre con me con grande affetto.

Un pensiero speciale va alla mia nonna. Non c'è stato esame che non iniziasse con una sua telefonata. Ogni volta mi accoglieva con il suo immancabile: ``Come sei bella, come sei intelligente'', e io rispondevo, ormai come un rito, con il nostro ``Cocoricò''. Un piccolo gesto, semplice ma speciale, che riusciva sempre a farmi sorridere e a regalarmi quella serenità di cui avevo bisogno prima di ogni prova. Ogni esame di questo percorso porta con sé anche un pezzetto di quelle telefonate e sono felice che, arrivata all'ultimo, possa finalmente dirti che sì, nonna: anche questa volta è andata.

A Simone Machetti, il mio ragazzo, grazie per aver condiviso con me ogni passo di questo percorso, dal primo all'ultimo giorno della magistrale. Grazie per avermi sostenuta nei momenti in cui tutto sembrava più difficile e per aver festeggiato con me ogni piccolo traguardo raggiunto. Grazie per avermi supportata e, qualche volta, anche sopportata, soprattutto durante le giornate in cui la tesi e gli esami sembravano occupare ogni mio pensiero.

Sei stato una presenza costante, capace di ricordarmi che un ostacolo alla volta si arriva sempre alla meta. Mi hai incoraggiata a credere di più in me stessa e nelle mie capacità, senza mai farmi sentire sola.

Ti ringrazio perché sei, ogni giorno, un esempio di impegno, determinazione e passione. La dedizione con cui affronti ciò che ami è per me uno stimolo continuo a dare il meglio di me stessa, a mettermi sempre in gioco e a non accontentarmi mai.

Avere una persona come te accanto ha reso questo percorso più leggero, più bello e ancora più significativo.

Alle mie amiche, in particolare Giorgia Ravetta, Alessia Vincenti, Martina Badano, Cecilia Coletti e Anna Trevisano, grazie per esserci sempre state. Grazie per aver compreso tutti quei ``non posso, devo studiare'', senza mai farmeli pesare, aspettando con pazienza il momento giusto per ritrovarci, ridere, raccontarci e condividere la vita. Sapere di poter sempre contare su di voi è stato un dono prezioso.

Infine, un grazie a me stessa. Per aver affrontato questo percorso con curiosità, determinazione e costanza. Per aver scelto di viverlo fino in fondo, accogliendone le soddisfazioni ma anche le difficoltà, senza smettere di mettermi in gioco.

Questo traguardo segna la conclusione di un capitolo importante della mia vita e, allo stesso tempo, l'inizio di nuove esperienze e nuovi obiettivi. Porto con me tutto ciò che questi anni mi hanno insegnato e mi auguro di affrontare ciò che verrà con la stessa curiosità, determinazione, costanza e gioia per la vita che mi hanno accompagnata fino a qui.

A questo traguardo e, mi auguro, ai tanti che verranno.

\setlength{\parskip}{0em}

\vspace{1.5cm}

\hfill \Candidato
